\section{Week 6: Efficiency, Scaling, and Mixture-of-Experts}
\label{sec:week6-moe}

\subsection{Learning Objectives}

By the end of this week, you should be able to:
\begin{itemize}
  \item Decompose an LLM workload along the three axes of compute,
        memory, and bandwidth, and identify which one is binding.
  \item Characterize data, tensor, pipeline, sequence (ring), and
        expert parallelism and pick a sensible composite for a given
        cluster.
  \item Derive the ZeRO staging (P/OS/G/activations) and relate it to
        the AdamW optimizer-state bookkeeping from Week~2.
  \item Write the MoE forward pass, including top-$k$ gating,
        load-balancing auxiliary loss, and token capacity factor.
  \item Contrast the Switch Transformer, GShard, and Mixtral
        architectural choices.
  \item Describe all-to-all communication in expert parallelism and
        its latency/bandwidth tradeoffs.
  \item Reason about the memory-vs-FLOPs regime shift in inference
        (MoE buys FLOPs but not memory).
  \item Enumerate MoE-specific failure modes (expert collapse, router
        instability, token drop, capacity starvation) and their
        operational signatures.
  \item Connect hardware (HBM vs SRAM, NVLink vs InfiniBand) to model
        design decisions.
\end{itemize}

\subsection{LLM Components Covered}

\textbf{Distributed training and serving}
\begin{itemize}
  \item Data, tensor, pipeline, sequence, expert parallelism.
  \item ZeRO and FSDP memory sharding.
  \item Activation recomputation and selective checkpointing.
\end{itemize}

\textbf{Sparse-activation architecture}
\begin{itemize}
  \item MoE FFN layers.
  \item Top-$k$ routing and load-balancing losses.
  \item Expert parallelism and all-to-all.
\end{itemize}

\textbf{Systems relevance}
\begin{itemize}
  \item Hardware-aware design (HBM, L2, NVLink, RDMA).
  \item Latency budgets and tail-latency guards.
  \item New failure modes from conditional compute.
\end{itemize}

\subsection{Conceptual Core: Three Axes of Scale}

A deployed LLM is constrained simultaneously by three resources:
\begin{center}
\begin{tabular}{|p{2.6cm}|p{3.4cm}|p{8.0cm}|}
\hline
Axis & Unit & What binds it \\
\hline
Compute   & FLOPs / s   & Forward + backward passes; matmul
                          throughput per GPU. \\
Memory    & bytes       & Weights + optimizer state + activations +
                          KV cache. \\
Bandwidth & bytes / s   & HBM-to-SRAM and inter-GPU interconnect;
                          often the true bottleneck. \\
\hline
\end{tabular}
\end{center}

Training a dense Transformer activates \emph{all} weights for every
token -- cost scales linearly with parameter count on the FLOPs axis
\emph{and} the memory axis \emph{and} the bandwidth axis.
Mixture-of-Experts (MoE) breaks this coupling:
\begin{quote}
MoE decouples \textbf{model capacity} (parameters that \emph{exist})
from \textbf{per-token compute} (parameters that \emph{activate}).
\end{quote}
A 1T-parameter MoE model with top-2 routing may activate only 30B
parameters per token, yielding the representational reach of a 1T
model at the inference FLOPs cost of a 30B model -- \emph{but with the
memory footprint of the 1T model}. This asymmetry drives most of what
is interesting about MoE at the systems level.

\subsection{Parallelism: The Five Axes}

See Figure~\ref{fig:week6-parallelism-axes}. In modern training
stacks these axes compose multiplicatively.

\begin{figure}[h]
\centering
\begin{tikzpicture}[
  ax/.style={draw, rounded corners, minimum width=3.5cm, minimum height=1.5cm,
             align=center, font=\footnotesize, fill=black!5},
  node distance=0.4cm and 0.4cm
]

\node[ax] (dp) {\textbf{Data parallel}\\full model on each GPU,\\different batches};
\node[ax, right=of dp] (tp) {\textbf{Tensor parallel}\\split matmuls across GPUs};
\node[ax, right=of tp] (pp) {\textbf{Pipeline parallel}\\split layers across GPUs};
\node[ax, below=of dp] (sp) {\textbf{Sequence / ring}\\split sequence axis across GPUs};
\node[ax, below=of tp] (ep) {\textbf{Expert parallel}\\distribute MoE experts};
\node[ax, below=of pp] (zero) {\textbf{ZeRO / FSDP}\\shard optimizer + params};

\end{tikzpicture}

\caption{Six parallelism dimensions layered in modern LLM training.
  Production stacks typically compose three or four of them
  simultaneously; expert parallelism applies only to MoE layers.}
\label{fig:week6-parallelism-axes}
\end{figure}

\subsubsection{Data Parallelism (DP)}

Each GPU holds a full copy of the model and processes a disjoint
microbatch. After the backward pass, gradients are averaged across
replicas via all-reduce. Scaling is near-linear on the compute axis
for batch sizes under the noise threshold of
\citet{mccandlish2018empirical}, and cheap to implement. Weakness: it
replicates parameters, so memory cost is $O(N)$ per replica and does
not benefit from cluster size.

\subsubsection{Tensor (Model) Parallelism (TP)}

Introduced at scale by Megatron-LM
\citep{shoeybi2019megatron,narayanan2021efficient}. Within a single
block, each linear layer $y = W x$ is split either row-wise or
column-wise across GPUs. For a column-wise split
$W = [W_1 \,|\, W_2]$,
\begin{equation}
y = W x = \begin{bmatrix} W_1 x \\ W_2 x \end{bmatrix},
\end{equation}
each GPU computes its slice locally and the results are concatenated
via an all-gather. TP requires fast intra-node interconnect (NVLink
or NVSwitch) because every forward and backward pass must exchange
activations; it is typically bounded to intra-node (8 GPUs on
current H100/H200 nodes).

\subsubsection{Pipeline Parallelism (PP)}

Different \emph{layers} live on different GPUs. A microbatch flows
forward through the stages, then backward.
GPipe~\citep{huang2019gpipe} and the 1F1B/interleaved schedules of
Megatron keep pipelines full by running many microbatches
concurrently at different stages. The residual idle time is the
\emph{pipeline bubble}, which scales as $(P - 1) / M$ for $P$ stages
and $M$ microbatches; large $M$ (i.e., large effective batch) shrinks
the bubble.

\subsubsection{Sequence / Ring Parallelism (SP)}

For long contexts, even a single sequence may exceed per-GPU
activation memory. Sequence
parallelism~\citep{li2023sequence,korthikanti2023reducing} shards
the sequence axis; \emph{Ring Attention}~\citep{liu2023ringattention}
arranges GPUs in a ring and rotates the $K$/$V$ blocks around the
ring so that every query sees every key without materializing the
full attention matrix anywhere. SP is what enables million-token
context training.

\subsubsection{Expert Parallelism (EP)}

Specific to MoE. Different \emph{experts} live on different GPUs.
When a token is routed to expert $i$, its activations must travel to
the GPU that holds $E_i$ -- this is the all-to-all communication
pattern (\S\,Expert Parallelism below).

\subsubsection{Composing the Axes}

A real 512-GPU Llama-3 training might decompose as $\mathrm{DP}=32$,
$\mathrm{TP}=4$, $\mathrm{PP}=4$, with $\mathrm{SP}=4$ inside TP for
long-context runs. The scheduling is typically expressed as a 3D
mesh and optimized by the framework (Megatron, DeepSpeed). The
engineering artifact that matters is the \emph{configuration file}
listing this decomposition, because it determines throughput, the
bubble size, and the blast radius of a single-GPU failure.

\subsection{ZeRO and FSDP}

Pure data parallelism replicates parameters, gradients, optimizer
state, and activations on every GPU -- memory is $O(N)$ per device
even on huge clusters. The Zero Redundancy Optimizer
(ZeRO)~\citep{rajbhandari2020zero} shards these tensors across data-
parallel ranks, reconstructing them only when needed:
\begin{itemize}
  \item \textbf{ZeRO-1}: shard the optimizer state (the $m_t, v_t$ of
        AdamW from Week~2). Memory per device drops by the DP degree
        with only an additional all-gather per step.
  \item \textbf{ZeRO-2}: additionally shard the gradients.
  \item \textbf{ZeRO-3}: additionally shard the parameters themselves;
        each forward pass all-gathers the needed shard, runs locally,
        then discards.
\end{itemize}
ZeRO-3 is memory-equivalent to tensor parallelism for the
parameter axis, but uses only DP-style all-gather communication
(typically over the higher-bandwidth fabric). PyTorch FSDP
\citep{zhao2023fsdp} is an open-source implementation of
ZeRO-3-style full sharding and has become the default for
open-source LLM training.

\subsubsection{Activation Recomputation}

Even after sharding parameters and optimizer state, activation memory
often dominates. Activation recomputation \citep{chen2016training}
(Week~2) drops intermediate activations and re-computes them in the
backward pass. \citet{korthikanti2023reducing} refined this to
\emph{selective} recomputation -- keep the activations for
compute-heavy blocks (attention, FFN) and recompute the cheap ones
(normalizations, reshapes) -- recovering most of the memory savings
at a $\sim 5\%$ compute overhead rather than the classical $\sim 33\%$.

\subsection{MoE Forward Pass}

\subsubsection{Layer Structure}

A dense Transformer block has an FFN:
$\mathrm{FFN}(x) = W_2\,\sigma(W_1 x)$. An MoE block replaces this
with an ensemble of experts $\{E_1, \ldots, E_N\}$ plus a router:
\begin{align}
\label{eq:week6-moe-router}
\mathbf{s}(x)        &= \mathrm{softmax}(W_r x)\in \Delta^{N-1}, \\
\label{eq:week6-moe-topk}
\mathcal{T}_k(x)     &= \mathrm{top}_k(\mathbf{s}(x)), \\
\label{eq:week6-moe-out}
y \;=\; \mathrm{MoE}(x)
  &= \sum_{i \in \mathcal{T}_k(x)} \tilde g_i(x)\,E_i(x),
\end{align}
with $W_r \in \mathbb{R}^{N \times d}$ the router projection and
$\tilde g_i(x) = s_i(x) / \sum_{j \in \mathcal{T}_k(x)} s_j(x)$ the
renormalized gate. Top-1 routing
\citep{fedus2021switch} activates one expert per token; top-2
\citep{shazeer2017outrageously, lepikhin2021gshard,jiang2024mixtral}
activates two. Figure \ref{fig:week6-moe-routing} diagrams the
operation.

\begin{figure}[h]
\centering
\begin{tikzpicture}[
  tok/.style={draw, rounded corners, minimum width=1.3cm, minimum height=0.7cm,
              align=center, font=\footnotesize},
  gate/.style={draw, rounded corners, minimum width=1.6cm, minimum height=0.8cm,
               align=center, font=\footnotesize, fill=black!5},
  expert/.style={draw, rounded corners, minimum width=1.5cm, minimum height=0.9cm,
                 align=center, font=\footnotesize},
  combine/.style={draw, rounded corners, minimum width=3.0cm, minimum height=0.8cm,
                  align=center, font=\footnotesize, fill=black!5},
  arrow/.style={->, thick, >=Stealth},
  dashedA/.style={->, thick, dashed, >=Stealth},
  node distance=0.55cm and 0.55cm
]

% Input tokens
\node[tok] (t1) {$x_1$};
\node[tok, right=of t1] (t2) {$x_2$};
\node[tok, right=of t2] (t3) {$x_3$};
\node[tok, right=of t3] (t4) {$x_4$};

% Router
\node[gate, below=0.55cm of t2.east, xshift=0.55cm] (router)
     {Router $g(\cdot)$};

\draw[arrow] (t1) -- (router);
\draw[arrow] (t2) -- (router);
\draw[arrow] (t3) -- (router);
\draw[arrow] (t4) -- (router);

% Experts
\node[expert, below=0.9cm of router, xshift=-3.3cm] (e1) {Expert\\$E_1$};
\node[expert, right=of e1] (e2) {Expert\\$E_2$};
\node[expert, right=of e2] (e3) {Expert\\$E_3$};
\node[expert, right=of e3] (e4) {Expert\\$E_4$};

% Top-2 routing: active (solid), inactive (dashed gray)
\draw[arrow]   (router) -- node[midway, left=-1pt, font=\scriptsize] {top-2} (e1);
\draw[arrow]   (router) -- (e2);
\draw[dashedA, draw=black!40] (router) -- (e3);
\draw[dashedA, draw=black!40] (router) -- (e4);

% Combine
\node[combine, below=0.8cm of e2.east, xshift=0.55cm] (comb) {$\sum_i g_i(x)\, E_i(x)$};

\draw[arrow] (e1) -- (comb);
\draw[arrow] (e2) -- (comb);
\draw[dashedA, draw=black!40] (e3) -- (comb);
\draw[dashedA, draw=black!40] (e4) -- (comb);

\node[tok, below=0.55cm of comb] (out) {$y$};
\draw[arrow] (comb) -- (out);

\node[right=0.7cm of router, font=\footnotesize\itshape, align=left]
     (note) {Per-token\\top-$k$ gate};

\end{tikzpicture}

\caption{MoE routing. Each token $x_t$ is independently scored by
  the router, the top-$k$ experts are selected, their outputs are
  combined by the (renormalized) gate weights, and the result is
  passed on. The dashed arrows represent inactive experts for this
  particular token.}
\label{fig:week6-moe-routing}
\end{figure}

\subsubsection{Capacity Factor}

In a distributed implementation, each expert has a fixed buffer size.
The \emph{capacity} per expert is
\begin{equation}
\label{eq:week6-capacity}
C = \lceil f \cdot k \cdot B / N \rceil,
\end{equation}
with $B$ the total tokens in the batch, $N$ the number of experts,
$k$ the top-$k$, and $f \ge 1$ the \emph{capacity factor}. If more
than $C$ tokens are routed to a given expert in a step, the overflow
tokens are either \emph{dropped} (their contribution is zero) or
routed to a backup expert. $f \in [1.0, 1.25]$ is typical. $f < 1$
saves compute but increases drop rate; large $f$ wastes compute on
empty slots.

\subsubsection{Load-Balancing Auxiliary Loss}

Routing is discrete and the router receives no direct gradient
through $\mathrm{top}_k$; without an auxiliary signal, the router can
collapse onto a few experts (``the rich get richer''). Switch
Transformer introduces
\begin{equation}
\label{eq:week6-load-balance}
\mathcal{L}_{\mathrm{aux}}
   \;=\; \alpha \cdot N \sum_{i=1}^{N} f_i \cdot P_i,
\end{equation}
where $f_i$ is the fraction of tokens dispatched to expert $i$ (a
discrete, non-differentiable measurement) and $P_i = \E[s_i(x)]$ is
the fraction of router \emph{probability mass} on expert $i$ (a
differentiable quantity). Minimizing $\mathcal{L}_{\mathrm{aux}}$
pushes $P_i$ down wherever $f_i$ is large -- if the router is already
sending many tokens to $E_i$, reduce your confidence in $E_i$.
Equilibrium is $f_i = P_i = 1/N$, i.e., uniform.

\subsubsection{\texorpdfstring{Router $z$-Loss}{Router z-Loss}}

\citet{zoph2022stmoe} introduced the router $z$-loss to stabilize
router logits:
\begin{equation}
\label{eq:week6-zloss}
\mathcal{L}_z \;=\;
  \beta \cdot \E_x\!\left[\bigl(\log\!\sum_i e^{(W_r x)_i}\bigr)^2\right],
\end{equation}
which penalizes the log-sum-exp of the router logits from growing
unboundedly. Without it, router logits can diverge in bf16 and
produce router NaNs in long training runs.

\subsection{Expert Parallelism and All-to-All}

In a distributed MoE, expert $i$ lives on GPU $\pi(i)$. When a token
is routed to $E_i$, its hidden state must \emph{travel} to GPU
$\pi(i)$. The communication pattern:
\begin{enumerate}
  \item \textbf{Dispatch all-to-all}: each GPU sends its outgoing
        tokens to the appropriate expert GPU.
  \item \textbf{Local expert compute}: each expert GPU applies its
        FFN to the tokens it received.
  \item \textbf{Combine all-to-all}: each GPU pulls back the results
        for its original tokens and weights them by the gate.
\end{enumerate}
All-to-all bandwidth scales as $O(B d)$ per step and saturates
inter-node interconnect quickly. MoE training is \emph{interconnect-
bound}, not compute-bound, on typical clusters; DeepSpeed-MoE
\citep{rajbhandari2022deepspeedmoe} and similar systems engineer
around this with expert-placement heuristics, communication
overlap, and hierarchical dispatch.

\subsection{MoE in Inference: Buys FLOPs, Not Memory}

At inference time, MoE still reduces FLOPs per token (only $k$
experts fire) but \emph{all experts must remain resident in GPU
memory} because the router can select any subset per token.
Consequences:
\begin{itemize}
  \item \textbf{Mixtral-8x7B} has 47B active parameters per token but
        a resident footprint of $\sim\!93$\,B parameters; it needs a
        cluster-memory profile similar to a 70B dense model.
  \item \textbf{Batch efficiency}: MoE benefits enormously from large
        batches because load-balancing statistics become tight; at
        batch size 1 the same expert can be loaded and unloaded from
        SRAM repeatedly.
  \item \textbf{Speculative decoding} (\citealt{leviathan2023speculative},
        Week~7) pairs well with MoE because the cheap draft model
        generates many candidates and the MoE verifies them in a
        large batch.
\end{itemize}

\subsection{Scaling Laws for MoE}

\citet{clark2022unified} fit a unified scaling law across dense and
routed models of the form
\begin{equation}
\label{eq:week6-unified-scaling}
L(N_{\mathrm{total}}, N_{\mathrm{active}}, D) \;\approx\;
  E + \frac{A}{N_{\mathrm{total}}^\alpha} +
      \frac{B}{N_{\mathrm{active}}^\gamma} +
      \frac{C}{D^\beta},
\end{equation}
with dense models recovered as $N_{\mathrm{total}} = N_{\mathrm{active}}$.
Two empirical takeaways:
\begin{itemize}
  \item At fixed training compute, an MoE with $N$ experts achieves
        lower loss than the dense model with the same per-token
        FLOPs.
  \item Scaling total parameters while holding active parameters
        fixed has diminishing returns; the ``sweet spot'' is
        typically 8--16 experts with top-2 routing.
\end{itemize}
These observations, combined with Chinchilla's FLOP accounting from
Week~2, drive the MoE proliferation of 2024--2025 (Mixtral, DeepSeek,
Qwen-MoE, Grok-1).

\subsection{MoE Failure Modes}

Conditional compute introduces failure modes invisible to dense
training.

\subsubsection{Expert Collapse}

The router concentrates all traffic on a small subset of experts
despite the auxiliary loss. Symptom: $f_i \to 0$ for most $i$, with a
few ``celebrity'' experts absorbing all tokens. Causes: too small
$\alpha$, router initialization, insufficient warmup on the aux loss,
distribution shift mid-training.

\subsubsection{Router Instability}

Router logits oscillate, sending a given token to expert 3 on one
step and expert 7 on the next -- nominally equivalent if $E_3$ and
$E_7$ are well-balanced, but in practice training slows to a crawl.
The $z$-loss (Eq.\ \eqref{eq:week6-zloss}) was designed to address
this; so was the switch to bf16 for router logits specifically.

\subsubsection{Token Drop}

When an expert's capacity $C$ (Eq.\ \eqref{eq:week6-capacity}) is
exceeded, overflow tokens are dropped. Small drop rates (few percent)
are benign; large drop rates starve the model of signal. Monitoring
the drop rate per layer is essential.

\subsubsection{Routing Brittleness at Inference}

A token that was reliably routed to $E_3$ during training may, in a
distribution-shifted inference context, route to an underfit expert.
Unlike dense models, MoE quality can vary across input distributions
in ways that are hard to anticipate.

\subsection{Hardware-Aware Design}

GPU memory has a hierarchy: SRAM (L1/L2/shared, $\sim\!20$\,MB, TB/s),
HBM (80\,GB, 3\,TB/s), NVLink/NVSwitch intra-node
($\sim\!900$\,GB/s), InfiniBand inter-node ($\sim\!400$\,GB/s). Cost
decreases, bandwidth plummets.

Design consequences:
\begin{itemize}
  \item \textbf{FlashAttention} (Week~4) exists because
        materializing $A \in \mathbb{R}^{T \times T}$ in HBM is a
        bandwidth tax. Its tiling moves the attention compute into
        SRAM.
  \item \textbf{Tensor parallelism} lives on NVLink; larger TP
        degrees work only within a node.
  \item \textbf{Pipeline and expert parallelism} tolerate InfiniBand
        because their communication is less frequent but larger.
  \item \textbf{Cluster topology} (rails, fat-tree, rail-optimized)
        dictates where to place which parallelism axis.
\end{itemize}

\subsubsection{The Arithmetic Intensity Lens}

A useful planning tool is \emph{arithmetic intensity}: FLOPs per byte
of memory traffic. An NVIDIA H100 has roughly $1000$ TFLOPs and
$3$ TB/s HBM, so its ``ridge point'' is near $330$ FLOPs/byte. A
matmul of size $M \times N \times K$ has arithmetic intensity
$\approx 2 M N K / (M K + N K + M N) \cdot b$ bytes$^{-1}$; for
typical LLM matmuls at training (large $M$ from batch) this is
hundreds of FLOPs/byte and compute-bound. Per-token inference of a
single prompt has $M = 1$ and is bandwidth-bound -- explaining why
\emph{batching} is the single largest lever for inference throughput
(Week~7).

\subsection{Systems Engineering Implications}

\begin{center}
\begin{tabular}{|p{3.4cm}|p{3.4cm}|p{6.8cm}|}
\hline
V-model stage & Efficiency / MoE property & System requirement or activity \\
\hline
Requirements & Inference latency budget & Prefer sparse if FLOPs
   dominate; dense if memory dominates. \\
 & Training compute budget & Chinchilla for dense, Eq.\,
   \eqref{eq:week6-unified-scaling} for MoE. \\
\hline
Architecture & 3D/4D parallelism mesh & Configure
   (DP, TP, PP, SP, EP) per workload. \\
 & Activation recomputation & Selective per Korthikanti et al. \\
 & ZeRO/FSDP stage & Choose by memory target and
   comm budget. \\
\hline
Implementation & Aux and $z$-losses & Part of training config;
   sweep $\alpha, \beta$. \\
 & Capacity factor & $f \in [1.0, 1.25]$ default;
   log drop rate. \\
\hline
V\&V & Per-expert token-count histograms
     & Alert on $\max f_i / \min f_i > 5$. \\
 & Routing regression tests & Compare routing decisions across
   model revisions on a fixed probe set. \\
\hline
Operations & Expert health dashboards & $f_i$, drop rate, router
   $z$-score, expert-specific loss. \\
 & Cluster topology awareness & Place EP on inter-node fabric; TP
   on intra-node NVLink. \\
\hline
\end{tabular}
\end{center}

\subsection{Human Factors and Failure Modes}

\begin{itemize}
  \item \textbf{MoE hype}. Do not assume MoE always wins; for
        memory-bound serving on a single-GPU budget, dense may be
        cheaper at fixed quality.
  \item \textbf{Behavior drift}. MoE outputs can shift subtly after
        a routing change; downstream evaluations should include
        representative distribution coverage.
  \item \textbf{Debugging overhead}. A failed MoE run can fail at any
        of: router, aux loss, all-to-all, capacity, expert placement.
        Dashboards and per-expert logs are non-optional.
  \item \textbf{Reproducibility}. MoE routing decisions depend on
        batch composition, so the same input can be processed
        differently depending on what's in the batch -- a form of
        non-determinism that surprises new operators.
\end{itemize}

\subsection{Bridges to Later Chapters}

\begin{itemize}
  \item \textbf{Week 7 (Adaptation, LoRA, Quantization)}: MoE
        interacts with quantization (expert-specific scales) and
        LoRA (per-expert adapters); serving tools (vLLM, TGI) need
        special support for MoE layers.
  \item \textbf{Week 8 (Alignment)}: instruction tuning of MoE models
        is sensitive to routing changes; aux-loss coefficients
        sometimes need re-tuning.
  \item \textbf{Week 9 (LLM Systems)}: batching strategies, prefix
        caching, and speculative decoding all interact with MoE
        sparsity.
  \item \textbf{Week 10 (Governance)}: conditional compute is a
        residual-risk item (routing under distribution shift); it
        should appear in the assurance case.
\end{itemize}

\subsection{Key References}

\begin{itemize}
  \item \citet{shazeer2017outrageously} --- original top-$k$
        sparsely-gated MoE layer.
  \item \citet{fedus2021switch} --- Switch Transformer (top-1
        routing, aux loss).
  \item \citet{lepikhin2021gshard} --- GShard and MoE sharding.
  \item \citet{jiang2024mixtral} --- Mixtral 8x7B.
  \item \citet{zoph2022stmoe} --- ST-MoE and router $z$-loss.
  \item \citet{clark2022unified} --- unified scaling laws for routed
        models.
  \item \citet{rajbhandari2022deepspeedmoe} --- DeepSpeed-MoE.
  \item \citet{rajbhandari2020zero}, \citet{zhao2023fsdp} --- ZeRO
        and FSDP.
  \item \citet{shoeybi2019megatron}, \citet{narayanan2021efficient}
        --- tensor/pipeline parallelism.
  \item \citet{huang2019gpipe} --- GPipe.
  \item \citet{korthikanti2023reducing} --- selective activation
        recomputation.
  \item \citet{li2023sequence}, \citet{liu2023ringattention} ---
        sequence and ring parallelism.
\end{itemize}

\subsection{Recommended University Lectures}

\begin{itemize}
  \item Stanford CS324 --- scaling, efficiency, and MoE,
        \citet{stanford_cs324}.
  \item Berkeley CS182 --- systems aspects of deep learning,
        \citet{berkeley_cs182}.
  \item MIT 6.S191 --- efficient deep learning systems,
        \citet{mit_6s191}.
\end{itemize}

\subsection{Practitioner Lab}

\begin{enumerate}
  \item \textbf{Arithmetic intensity audit}. For one forward pass of
        a reference LLM at batch sizes $1, 8, 64$, compute
        arithmetic intensity per matmul. Plot against the ridge
        point of your GPU.
  \item \textbf{Data-parallel scaling study}. Train a small model on
        1, 2, 4, 8 GPUs. Plot throughput vs GPU count; identify
        where the comm-bound regime starts.
  \item \textbf{ZeRO stages}. Run ZeRO-1, -2, -3 on the same model
        and batch. Report peak memory per device and wall clock per
        step.
  \item \textbf{MoE from scratch}. Implement Eq.\
        \eqref{eq:week6-moe-router}--\eqref{eq:week6-moe-out} for a
        small Transformer. Add Eq.\ \eqref{eq:week6-load-balance}.
        Train on a toy task.
  \item \textbf{Capacity factor sweep}. Sweep $f \in \{0.8, 1.0,
        1.25, 2.0\}$. Log drop rate, load imbalance, and final loss.
  \item \textbf{Expert-collapse triggering}. Set $\alpha = 0$;
        confirm the router collapses onto few experts within a few
        thousand steps.
  \item \textbf{Router $z$-loss stability}. Train with and without
        Eq.\ \eqref{eq:week6-zloss}; log the max absolute router
        logit per step.
  \item \textbf{MoE serving memory}. Profile peak resident memory
        for Mixtral-8x7B at batch 1 and batch 32; reconcile with the
        total-parameter count.
\end{enumerate}

\subsection{Week 6 Takeaway}

\begin{quote}
Scaling LLMs is no longer a modeling problem alone -- it is a
distributed systems problem. Mixture-of-Experts decouples model
capacity from per-token compute and so is central to frontier
training, but it imports every pathology of conditional compute:
expert collapse, router instability, token drop, and interconnect-
bound training. The engineering success of MoE is not a question of
mathematical elegance; it is a question of whether the cluster's
topology, the aux losses, and the capacity factors have been tuned
jointly enough to make the sparsity actually pay.
\end{quote}
